\subsection{Requirements}
Based on the background and state of art section, a couple of requirements were made. These outlined the product and the accompanying system. Afterwards, in the problem analysis section, the final requirements were made. Those requirements drew from the UML diagrams, such as context and deployment diagram, which was used to create an overview of the system, its ability, the hardware parts and how it all fit together. 

\subsection{System architecture}
In terms of system architecture, a choice had to be made in terms of hardware, what components do we need, and what form-factor do we want to use. Based on the collective knowledge and idea of what the product should be able to do, there was a need for the micro-controller to have Bluetooth and an accelerometer, and preferably some easy way to integrate a battery to allow for wireless operation. Furthermore, there was a wish for an all-in-one product, so that the final product would be compact, and easier to develop. The need for it to be small was because the vision of the product being a wristband was already on the table. This is why a micro-controller ESP32C6 dev board made specifically for smart watches was chosen. This, along with a 3.7-volt battery, concludes the hardware. For the software, the Arduino IDE was used along with C++ code for the microcontroller, whereas for the application, VS Code with Typescript using the React framework was used. 

\subsection{Testing}
To ensure that the system functioned in accordance with the specified requirements, a comprehensive evaluation strategy was established. The wristband’s battery performance was assessed through multiple voltage measurements. The algorithm’s performance was evaluated using standard performance metrics, including sensitivity, specificity, and overall accuracy.

The application code was verified through a suite of unit tests designed to validate individual functions and features. In addition, a user evaluation was conducted to assess the intuitiveness and usability of both the wristband and the application. Finally, an integration evaluation was performed to examine the application as a complete system, simulating real user interaction.
